\section{Continuous rational norm pushforward}

Let
\[
\Brat=\ell^1(\{p:p\text{ rational prime}\},\log p).
\tag{5.1}
\]
On tagged basis vectors define
\[
\nu[\gamma,n,\mathfrak q]=f_{\mathfrak q}[p],
\qquad \mathfrak q\mid p.
\tag{5.2}
\]

\begin{proposition}[Norm-one extension]\label{prop:pushforward}
The map $\nu$ extends uniquely to a bounded linear map
$\nu:\Btag\to\Brat$ with $\|\nu\|=1$.  For every packet,
\[
\|\nu D_{\gamma,n}\|_{\rm rat}
=\|D_{\gamma,n}\|_{\rm tag}.
\tag{5.3}
\]
Consequently $\nu\mathcal G(s,u)$ is obtained by termwise pushforward in the
domain of Theorem \ref{thm:germ}.
\end{proposition}

\begin{proof}
For finitely supported complex coefficients,
\[
\begin{aligned}
\|\nu c\|_{\rm rat}
&=\sum_p\left|
\sum_{\substack{\gamma,n,\mathfrak q\\\mathfrak q\mid p}}
f_{\mathfrak q}c_{\gamma,n,\mathfrak q}\right|\log p\\
&\le\sum_{\gamma,n,\mathfrak q}
|c_{\gamma,n,\mathfrak q}|f_{\mathfrak q}\log p
=\|c\|_{\rm tag}.
\end{aligned}
\]
Thus $\nu$ extends with norm at most one.  Equality on every basis vector
shows that the norm is one.  Packet coefficients are nonnegative, so the
triangle inequality is an equality and (5.3) follows.  Continuity and the
normal convergence of Theorem \ref{thm:germ} justify termwise application.
\end{proof}

The proposition solves the topology problem without solving the information
problem.  The finite tagged certificate already has nontrivial kernel after
rational pushforward: several orbit/index/prime-ideal atoms can map to the
same rational-prime basis vector.  Therefore (5.3) is not an injectivity
statement and does not define a rational-prime clock orbit by orbit.
